\chapter{Data Collection}
\label{ch:collection}

Two outdoor recording campaigns were carried out with Ricbot during this
work (\Cref{tab:campaigns}).
The earlier \emph{Initial Development} campaign was recorded on and around
the DTU Lyngby campus between February and March 2026 and was used to
design and validate the processing pipeline itself.
The later \emph{Main Agricultural} campaign, recorded at a working farm
in April 2026, forms the primary dataset on which all classifiers in this
thesis are trained, tuned, and cross-validated.

\begin{table}[H]
\centering
\caption{Overview of the two data-collection campaigns.}
\label{tab:campaigns}
\begin{tabular}{lllp{4.4cm}}
\toprule
\textbf{Campaign} & \textbf{Period} & \textbf{Environment} & \textbf{Role} \\
\midrule
Initial Development  & Feb–Mar 2026  & Campus / urban paths
               & Pipeline development and cross-environment test \\
Main Agricultural    & Apr 2026      & Working farm
               & Primary training, tuning, and evaluation \\
\bottomrule
\end{tabular}
\end{table}

\section{Initial Development Dataset}
\label{sec:initial_dev_dataset}

The Initial Development dataset (stored under \texttt{data/raw/Main}
in the project repository) was recorded across four sessions on and
around the DTU Lyngby campus between February and March 2026.
The intent of this campaign was twofold: to gather a first body of
labelled IMU and odometry data covering a wide range of outdoor
surfaces, and to use that data to make the practical pipeline-design
decisions that the rest of the thesis depends on — in particular,
the analysis-window length, the minimum amount of useful data required
per run for stable Leave-One-Run-Out (LORO) cross-validation, and the
approximate target number of runs needed per terrain class before
results become reproducible.
For these reasons it is used in the final experiments as an independent
\emph{test set} for cross-environment generalisation rather than as part
of the main training pool.

The four campus sessions are summarised in \Cref{tab:main_sessions}.
Each session followed a predefined route covering several terrain types
in sequence, and the robot was teleoperated at approximately constant
speed within each terrain segment whenever feasible.
Weather conditions varied substantially across the campaign, ranging
from wet snow to dry pavement, which is one of the reasons this dataset
was retained as a stress-test rather than promoted to the primary
training set.

\begin{table}[H]
\centering
\caption{Sessions in the Initial Development dataset.}
\label{tab:main_sessions}
\begin{tabular}{llp{5.0cm}p{3.4cm}}
\toprule
\textbf{Run ID} & \textbf{Date} & \textbf{Conditions} & \textbf{Terrains present} \\
\midrule
\texttt{log\_20260223\_142511} & 23 Feb 2026 & Snowy and wet; wet snow on grass &
   grass, smooth, muddy dirt track \\
\texttt{log\_20260226\_102148} & 26 Feb 2026 & Less snowy, mostly wet &
   grass, smooth, muddy dirt track \\
\texttt{log\_20260309\_141435} &  9 Mar 2026 & Mostly dry; rough sections &
   grass, smooth, dry dirt track, cobblestone \\
\texttt{log\_20260326\_120021} & 26 Mar 2026 & Mostly dry &
   grass, smooth, cobblestone \\
\bottomrule
\end{tabular}
\end{table}

\placeholder{Optional: add total recording duration per session and the
total labelled duration per terrain class (in minutes). These numbers
can be filled in from the per-run QC reports in \texttt{reports/raw/}.}

\section{Main Agricultural Dataset}
\label{sec:farm_dataset}

The Main Agricultural dataset (stored under \texttt{data/raw/Farm} in the
project repository) was recorded on two field days at a working farm
during April 2026 and constitutes the primary dataset for this thesis.
It was collected after the Initial Development campaign once the
preprocessing pipeline, the windowing length, and the LORO cross-validation
protocol had been fixed; data was therefore gathered with a specific target
number of runs per class in mind, on a single site, under broadly
consistent weather conditions, so that the resulting class distributions
are balanced enough for stable cross-validated learning.

Five runs were recorded in total, spread across the two field days,
following an identical route in each case (\Cref{tab:farm_sessions}).
The route was selected so that each of the three target terrain classes
appears multiple times per run with comparable segment durations.
All five runs were carried out under dry, sunny conditions with light
wind, which keeps environmental confounds approximately constant across
runs and makes the variation between runs predominantly attributable to
operator-driven differences (path, speed) rather than weather.

\begin{table}[H]
\centering
\caption{Sessions in the Main Agricultural dataset.}
\label{tab:farm_sessions}
\begin{tabular}{lllp{4.0cm}}
\toprule
\textbf{Run ID} & \textbf{Date} & \textbf{Time}      & \textbf{Conditions} \\
\midrule
\texttt{Run1} & 21 Apr 2026 & 13:34–13:54 & Sunny, slightly windy, dry \\
\texttt{Run2} & 21 Apr 2026 & 13:56–14:17 & Sunny, slightly windy, dry \\
\texttt{Run3} & 28 Apr 2026 & 13:30–13:55 & Sunny, windy, dry \\
\texttt{Run4} & 28 Apr 2026 & 13:43–13:56 & Sunny, slightly windy, dry \\
\texttt{Run5} & 28 Apr 2026 & 13:56–14:09 & Sunny, slightly windy, dry \\
\bottomrule
\end{tabular}
\end{table}

\placeholder{Optional figure: aerial photograph or simple sketch of the
farm site with the route overlaid and the three terrain types coloured.}

\section{Terrain Classes and Labelling Strategy}
\label{sec:terrain_classes_and_labels}

\subsection{Primary Terrain Classes (Farm)}

Three terrain classes are considered in the primary classification problem,
all recorded at the agricultural site during the Main Agricultural campaign
(\Cref{tab:farm_terrain_classes}).

\begin{table}[H]
\centering
\caption{Primary terrain classes (Main Agricultural dataset).}
\label{tab:farm_terrain_classes}
\begin{tabular}{lp{9cm}}
\toprule
\textbf{Class} & \textbf{Description} \\
\midrule
\texttt{grass}       & Short-to-medium grass on firm soil between crop rows.
                       Moderate lateral and vertical vibration from
                       intermittent blade and tussock contact. \\
\texttt{dirt\_track} & Compacted unpaved farm track.
                       Higher-frequency surface roughness with moderate
                       vertical amplitude. \\
\texttt{soil}        & Tilled or bare agricultural soil.
                       Softer than the dirt track, with lower-frequency
                       deflection of the chassis and increased tyre slip. \\
\bottomrule
\end{tabular}
\end{table}

\begin{figure}[H]
    \centering
    \placeholder{Three small panels — one representative photograph of
    grass, dirt track, and soil from the farm site — captioned with the
    class name.}
    \caption{Representative photographs of the three primary terrain
    classes.}
    \label{fig:farm_terrain_photos}
\end{figure}

\subsection{Extended Terrain Classes (Initial Development)}

The Initial Development campaign covers a wider variety of urban surfaces
and is used in this thesis to study how the models trained on the Main
Agricultural data behave when applied to a clearly different environment.
Five classes are present in the campus logs
(\Cref{tab:main_terrain_classes}).

\begin{table}[H]
\centering
\caption{Terrain classes present in the Initial Development dataset.}
\label{tab:main_terrain_classes}
\begin{tabular}{lp{9cm}}
\toprule
\textbf{Class} & \textbf{Description} \\
\midrule
\texttt{smooth\_terrain}    & Hard, flat surface (asphalt road or
                              compacted concrete).
                              Low vibration amplitude across all axes. \\
\texttt{grass}              & Short-to-medium grass on firm soil.
                              Moderate lateral and vertical vibration. \\
\texttt{dry\_dirt\_track}   & Compacted unpaved path under dry conditions. \\
\texttt{muddy\_dirt\_track} & Unpaved path with surface moisture or wet snow;
                              increased tyre slip, lower dominant frequency. \\
\texttt{cobblestone}        & Irregular stone paving; high-amplitude,
                              broadband vibration with prominent impulsive
                              events. \\
\bottomrule
\end{tabular}
\end{table}

\begin{figure}[H]
    \centering
    \placeholder{Five small panels — one representative photograph per
    Initial Development terrain class.}
    \caption{Representative photographs of the five Initial Development
    terrain classes.}
    \label{fig:main_terrain_photos}
\end{figure}

\placeholder{Decide whether to keep the Initial Development 5-class
breakdown as-is when comparing against the Farm-trained models, or to
collapse it into a 3-class scheme (\emph{smooth}, \emph{grass},
\emph{rough}) so that the Initial Development data can be evaluated
directly against models trained on the Farm 3-class scheme. The
\texttt{*\_3class} notebooks (07a / 07b / 08b / 09b / 10b) suggest the
latter; if you commit to that, add a short paragraph here explaining
the merge and the reasoning.}

\subsection{Labelling Procedure}

All terrain labels were assigned manually after each session.
During each run, a paper log book was used to note the approximate Unix
time at which the robot crossed a transition between terrain types.
After the run, those times were refined by replaying the synchronised
sensor streams and the pose track, and the final time ranges were stored
as a per-run JSON configuration file (\texttt{labels\_config.json}).
\Cref{lst:labels_json} gives a representative example from
\texttt{Farm/Run1}.

\begin{lstlisting}[language=, caption={Example label configuration (\texttt{labels\_config.json}).}, label={lst:labels_json}]
{
  "grass":      [[1776764662, 1776764730],
                 [1776764798, 1776764844]],
  "dirt_track": [[1776764860, 1776765022]],
  "soil":       [[1776764533, 1776764573],
                 [1776764590, 1776764643]]
}
\end{lstlisting}

Time ranges are absolute Unix epoch seconds matching the timestamp column
of the raw log files.
Each $[t_{\text{start}}, t_{\text{end}}]$ pair defines a contiguous
interval during which the robot was confidently driving on a single
terrain type; unlabelled regions (turns, robot stops, ambiguous
transitions, or short periods in which the surface could not be
identified visually) are discarded during preprocessing rather than
being assigned to a catch-all class.

This labelling scheme is intentionally conservative: it sacrifices a
fraction of the recorded data in exchange for cleaner per-window class
assignments downstream.
Because every window inherits its label from the time range of its
mid-point, conservative ranges also prevent transition samples from
being silently mixed into a single class --- the dynamics of explicit
terrain transitions are studied separately in \Cref{ch:transitions}.